%% chapters/02_literature_review.tex

\section{Literature Review}
\label{sec:litreview}

\subsection{Class Imbalance in Fraud Detection}

The severe class imbalance in fraud datasets (typically 0.1--5\% positives)
renders standard accuracy misleading and challenges gradient-based learners
whose loss functions treat all samples equally.
\citet{chawla2002smote} introduced SMOTE, which generates synthetic minority
samples by interpolating between existing minority instances in feature space,
providing a robust alternative to naive oversampling.
\citet{han2005borderline} extended this with Borderline-SMOTE, focusing
synthetic generation on the decision boundary where the model's uncertainty
is greatest.
We compare five imbalance strategies in Section~\ref{sec:methodology}.

\subsection{Gradient Boosting for Fraud Detection}

\citet{chen2016xgboost} introduced XGBoost, which combines regularised
gradient boosting with an efficient tree construction algorithm (\textit{hist}
method) and built-in support for class weighting via the
\texttt{scale\_pos\_weight} hyperparameter.
XGBoost has consistently placed among the top performers in fraud detection
benchmarks on Kaggle, including the IEEE-CIS competition.
We use Optuna \citep{akiba2019optuna} with the Tree-structured Parzen Estimator
(TPE) sampler for Bayesian hyperparameter optimisation over 100 trials.

\subsection{Unsupervised Anomaly Detection}

\citet{liu2008isolation} proposed Isolation Forest, which exploits the
intuition that anomalies are few and different, isolating them with
shorter average path lengths in random trees.
Unlike density-based methods, Isolation Forest scales linearly with dataset
size and handles high-dimensional data without requiring a distance metric.

\citet{pimentel2014review} provide a comprehensive taxonomy of novelty
detection methods. Autoencoder-based anomaly detection, a neural approach,
trains a bottleneck network to reconstruct normal data; anomalies
produce high reconstruction error because they lie off the learned normal
manifold \citep{pimentel2014review}.
We implement a PyTorch Autoencoder trained exclusively on legitimate
transactions, using the Adam optimiser \citep{kingma2015adam} with
a ReduceLROnPlateau scheduler and gradient norm clipping.

\subsection{Evaluation Metrics for Imbalanced Data}

\citet{davis2006precision} demonstrate that the Precision-Recall (PR) curve
provides a more informative picture than the ROC curve when classes are
highly imbalanced, because it directly characterises the trade-off between
false alarms and missed detections.
\citet{chicco2020mcc} show that the Matthews Correlation Coefficient (MCC)
is a more balanced single-number summary than F1 for imbalanced binary
classification, as it accounts for all four cells of the confusion matrix.
We report AUPRC as the primary metric, with AUROC, MCC, F1, and Brier Score
as secondary metrics.

\subsection{Cost-Sensitive Learning and Threshold Selection}

\citet{elkan2001cost} formalise the cost-sensitive learning framework,
showing that the optimal threshold for a calibrated classifier can be
derived analytically from the misclassification cost ratio:
\begin{equation}
    \theta^* = \frac{c_{FP}}{c_{FP} + c_{FN}}
    \label{eq:elkan_threshold}
\end{equation}
In practice, because classifiers are not perfectly calibrated, we sweep
the threshold numerically and select $\theta^*$ by minimising the Expected
Cost curve over the validation set.

\subsection{Explainability via SHAP}

\citet{lundberg2017shap} introduced SHAP (SHapley Additive exPlanations),
grounding feature attribution in cooperative game theory.
TreeExplainer provides exact SHAP values for tree-based models in
polynomial time; KernelExplainer provides model-agnostic approximations
via background sampling.
We use TreeExplainer for XGBoost and KernelExplainer for Logistic Regression,
embedding the top-3 SHAP contributors in every API response.

\subsection{Target Encoding}

\citet{micci2001target} introduced smoothed target encoding, replacing
categorical levels with their smoothed empirical fraud rate.
To prevent target leakage, we apply cross-fitting: each fold's
encoding is computed from the out-of-fold training labels,
ensuring no transaction sees its own label during feature computation.
